% Сборка: lualatex lab1-report.tex
\documentclass[11pt,a4paper]{article}
\usepackage[a4paper,margin=2cm,top=1.6cm,bottom=1.8cm]{geometry}
\usepackage{fontspec}
\usepackage[russian]{babel}
\usepackage{amsmath}
\usepackage{unicode-math}
\setmainfont{Liberation Serif}
\setmonofont{Liberation Mono}[Scale=MatchLowercase]
\setmathfont{TeX Gyre Termes Math}
\usepackage{titlesec}
\usepackage{enumitem}
\usepackage{microtype}

\titleformat{\section}{\large\bfseries}{\thesection.}{0.5em}{}
\titlespacing*{\section}{0pt}{1.0ex plus .3ex}{0.5ex}
\setlength{\parindent}{0pt}
\setlength{\parskip}{0.45em}
\pagestyle{plain}

\usepackage[hidelinks,pdftitle={Лабораторная работа № 1. Цветовые модели},pdfauthor={Aliaksei Charnukha}]{hyperref}

\newcommand{\xmax}{x_{\max}}
\newcommand{\xmin}{x_{\min}}

\begin{document}

\begin{center}
  {\LARGE\bfseries Лабораторная работа № 1}\\[0.4em]
  {\large Цветовые модели. Чётный вариант: CMYK — RGB — HSV}
\end{center}

\section{Цель}
Реализовать интерактивное преобразование цвета между моделями CMYK, RGB и HSV.

\section{Основные формулы}
RGB задаётся в диапазоне $0\text{–}255$, CMYK, $S$ и $V$ — в процентах
$0\text{–}100$, $H$ — в градусах $0\text{–}360$. В формулах ниже
$r, g, b, c, m, y, k, s, v$ нормированы к отрезку $[0, 1]$:
\[
  r = \frac{R}{255}, \qquad g = \frac{G}{255}, \qquad b = \frac{B}{255}.
\]

\textbf{RGB $\to$ CMYK.} $\xmax = \max(r, g, b)$:
\[
  k = 1 - \xmax, \qquad
  c = \frac{\xmax - r}{\xmax}, \qquad
  m = \frac{\xmax - g}{\xmax}, \qquad
  y = \frac{\xmax - b}{\xmax}.
\]
При $\xmax = 0$ полагаем $c = m = y = 0$, $k = 1$.

\textbf{CMYK $\to$ RGB.}
\[
  R = 255\,(1 - c)(1 - k), \qquad
  G = 255\,(1 - m)(1 - k), \qquad
  B = 255\,(1 - y)(1 - k).
\]

\textbf{RGB $\to$ HSV.} $\xmax = \max(r, g, b)$, $\xmin = \min(r, g, b)$,
$\Delta = \xmax - \xmin$:
\[
  v = \xmax, \qquad
  s = \begin{cases} 0, & \xmax = 0,\\[2pt] \dfrac{\Delta}{\xmax}, & \xmax > 0, \end{cases}
  \qquad
  H = \begin{cases}
    60 \cdot \dfrac{g - b}{\Delta}, & \xmax = r,\\[8pt]
    60 \cdot \left(\dfrac{b - r}{\Delta} + 2\right), & \xmax = g,\\[8pt]
    60 \cdot \left(\dfrac{r - g}{\Delta} + 4\right), & \xmax = b,
  \end{cases}
  \quad (\Delta \neq 0).
\]
$H$ приводится к промежутку $[0, 360)$. При $\Delta = 0$ тон не определён.

\textbf{HSV $\to$ RGB.}
\[
  C = v \cdot s, \qquad
  X = C \left(1 - \left| \left(\frac{H}{60} \bmod 2\right) - 1 \right|\right), \qquad
  d = v - C.
\]
Для шести последовательных секторов по $60^\circ$ выбирается тройка $(r', g', b')$:
\[
  (C, X, 0),\; (X, C, 0),\; (0, C, X),\; (0, X, C),\; (X, 0, C),\; (C, 0, X),
\]
после чего $(R, G, B) = 255 \cdot (r' + d,\; g' + d,\; b' + d)$.

\section{Ход работы}
Приложение написано на HTML, CSS и JavaScript. Изменение CMYK или HSV
переводится в RGB, из которого пересчитываются остальные модели.

Координаты задаются числовыми полями, ползунками и системной палитрой, для HSV
есть цветовой квадрат с полосой тона. Цвет показывается образцом и HEX-кодом.
Некорректный ввод игнорируется, значения округляются только при выводе.

Формулы вынесены в отдельный модуль и проверены юнит-тестами.

\section{Результаты}
\begin{itemize}[label=--, leftmargin=1.5em, itemsep=0.2em, topsep=0.2em]
  \item Разработано веб-приложение для пересчёта цвета между моделями CMYK,
        RGB и HSV; изменение любой координаты сразу пересчитывает две другие модели.
  \item Реализованы три способа ввода: числовые поля, ползунки и системная
        палитра; для HSV добавлен цветовой квадрат с полосой тона.
  \item Формулы вынесены в отдельный модуль и проверены 40 юнит-тестами, в том
        числе на 529 эталонах стандартной библиотеки Python \texttt{colorsys}.
\end{itemize}

\section{Выводы}
\begin{itemize}[label=--, leftmargin=1.5em, itemsep=0.2em, topsep=0.2em]
  \item RGB удобно использовать как промежуточную модель: переход
        CMYK $\leftrightarrow$ HSV сводится к двум преобразованиям через RGB.
  \item Переход RGB $\to$ CMYK неоднозначен: один цвет можно получить разными
        сочетаниями красок. Выбирается вариант с максимальным $K$. Учебная
        формула без ICC-профиля не описывает реальную печать.
  \item У серых цветов ($\Delta = 0$) тон $H$ не определён, поэтому при
        пересчёте сохраняется предыдущий тон, и ползунок $H$ не сбрасывается.
  \item HSV удобнее для выбора цвета человеком: тон, насыщенность и яркость
        меняются независимо, поэтому цвет проще подобрать мышью на цветовом
        квадрате, чем числами RGB.
\end{itemize}

\end{document}
